% VI/38 detailed challenge solutions -- dossier-preserving refinement.

% VI38-DETAILED-CHALLENGE-01
\subsection*{Challenge VI.38.C1 solution}
\begin{solution}
Let
\[
S=\bigoplus_{r\ge0}S_r
\]
be finitely generated standard graded and let
\[
I=(F_0,\ldots,F_m)
\]
with all \(F_i\) homogeneous of the same positive degree \(d\).  Regrading by the
\(d\)-th Veronese gives the degree-preserving map
\[
k[y_0,\ldots,y_m]\longrightarrow S^{(d)},
\qquad y_i\longmapsto F_i.
\]
Its projective domain is the complement of the base locus
\[
V_+(I)=\{\mathfrak p\in\Proj S:I\subseteq\mathfrak p\}.
\]
Therefore the associated map is everywhere defined exactly when
\[
V_+(I)=\varnothing.
\]

For a finitely generated standard graded ring this has the intrinsic radical
criterion
\[
\boxed{V_+(I)=\varnothing
\iff
S_+\subseteq\sqrt I}.
\]
Indeed, if \(S_+\subseteq\sqrt I\), then any homogeneous prime containing \(I\)
contains \(S_+\), so it cannot represent a point of \(\Proj S\).

Conversely, suppose \(V_+(I)=\varnothing\). Choose degree-one generators
\(x_0,\ldots,x_r\) of \(S_+\). The standard opens \(D_+(x_j)\) cover \(\Proj S\).
Since \(V_+(I)\cap D_+(x_j)=\varnothing\), the localized ideal generated by the
\(F_i\) is the unit ideal in every projective chart. Equivalently, some power of
\(x_j\) lies in \(I\) after clearing homogeneous denominators; hence
\[
x_j\in\sqrt I
\]
for every \(j\). Thus \(S_+\subseteq\sqrt I\).

Because \(S_+\) is finitely generated, the radical containment is also equivalent
to a power containment
\[
\boxed{S_+^N\subseteq I\quad\text{for some }N\gg0}.
\]
Thus the universal base-locus test can be stated without reference to coordinates
of individual points:
\[
\boxed{
[F_0:\cdots:F_m]\text{ is defined on all of }\Proj S
\iff S_+\subseteq\sqrt{(F_0,\ldots,F_m)}.
}
\]
\end{solution}

% VI38-DETAILED-CHALLENGE-02
\subsection*{Challenge VI.38.C2 solution}
\begin{solution}
Write the degree-\(d\) Veronese map of \(\mathbb P^1\) as
\[
\nu_d([s:t])
=[s^d:s^{d-1}t:\cdots:st^{d-1}:t^d]
=[Z_0:Z_1:\cdots:Z_d].
\]
Consider the \(2\times d\) catalecticant matrix
\[
C=
\begin{pmatrix}
Z_0&Z_1&\cdots&Z_{d-1}\\
Z_1&Z_2&\cdots&Z_d
\end{pmatrix}.
\]
Let \(J\) be the ideal generated by its \(2\times2\) minors
\[
Z_iZ_{j+1}-Z_{i+1}Z_j,
\qquad 0\le i<j\le d-1.
\]
Under the coordinate-ring map
\[
\theta:k[Z_0,\ldots,Z_d]\longrightarrow k[s,t]^{(d)},
\qquad
Z_i\longmapsto s^{d-i}t^i,
\]
each such minor maps to zero. Thus
\[
J\subseteq\ker\theta.
\]

We prove the reverse inclusion by a straightening argument.  A monomial
\[
M=Z_{i_1}\cdots Z_{i_r}
\]
maps to
\[
s^{rd-(i_1+\cdots+i_r)}t^{i_1+\cdots+i_r}.
\]
Hence its image is determined by its degree \(r\) and its index sum
\[
q=i_1+\cdots+i_r.
\]
The minor relations give
\[
Z_iZ_j\equiv Z_{i+1}Z_{j-1}\pmod J
\qquad(j-i\ge2).
\]
Repeatedly applying this operation decreases the spread of the indices while
preserving their number and sum. Therefore every degree-\(r\) monomial is congruent
to a balanced normal form
\[
Z_a^{\,r-b}Z_{a+1}^{\,b},
\qquad
q=ar+b,
\quad 0\le b<r,
\]
with the endpoint convention when \(a=d\).  This normal form is uniquely determined
by \((r,q)\).

Consequently any two monomials with the same image under \(\theta\) are congruent
modulo \(J\).  Group a kernel polynomial by image monomial.  In each group the
coefficients sum to zero, so that group is a linear combination of differences of
monomials with the same image; each such difference lies in \(J\). Hence
\[
\boxed{\ker\theta=J}.
\]
Thus the homogeneous ideal of the rational normal curve is generated by the
\(2\times2\) minors of the catalecticant matrix \(C\).
\end{solution}

% VI38-DETAILED-CHALLENGE-03
\subsection*{Challenge VI.38.C3 solution}
\begin{solution}
Let
\[
p_1:\mathbb P^m\times\mathbb P^n\to\mathbb P^m,
\qquad
p_2:\mathbb P^m\times\mathbb P^n\to\mathbb P^n
\]
be the projections, and let \(a,b\ge1\). First apply the Veronese embeddings
\[
\nu_a:\mathbb P^m\hookrightarrow\mathbb P^{M},
\qquad
M=\binom{m+a}{a}-1,
\]
\[
\nu_b:\mathbb P^n\hookrightarrow\mathbb P^{N},
\qquad
N=\binom{n+b}{b}-1.
\]
Their product is a closed immersion
\[
\mathbb P^m\times\mathbb P^n
\hookrightarrow
\mathbb P^M\times\mathbb P^N.
\]
Compose with the Segre closed immersion
\[
\mathbb P^M\times\mathbb P^N
\hookrightarrow
\mathbb P^{(M+1)(N+1)-1}.
\]
The resulting Segre--Veronese map is therefore a closed immersion.

Its homogeneous coordinates are all products
\[
x^\alpha y^\beta,
\qquad
|\alpha|=a,
\quad
|\beta|=b,
\]
so they are exactly the bihomogeneous monomials of bidegree \((a,b)\). The target
dimension is
\[
\boxed{
\binom{m+a}{a}\binom{n+b}{b}-1
}.
\]

The line bundle pulled back from \(\OO(1)\) on the final projective space is
computed in two steps. The Segre embedding pulls back \(\OO(1)\) to
\[
\OO_{\mathbb P^M\times\mathbb P^N}(1,1)
=p_1^*\OO_{\mathbb P^M}(1)\otimes p_2^*\OO_{\mathbb P^N}(1).
\]
Each Veronese pullback satisfies
\[
\nu_a^*\OO(1)\cong\OO(a),
\qquad
\nu_b^*\OO(1)\cong\OO(b).
\]
Hence the composite pulls back \(\OO(1)\) to
\[
\boxed{
\OO(a,b)
:=p_1^*\OO_{\mathbb P^m}(a)\otimes p_2^*\OO_{\mathbb P^n}(b).
}
\]
Since the composite is a closed immersion, \(\OO(a,b)\) is very ample for
\(a,b\ge1\).
\end{solution}

% VI38-DETAILED-CHALLENGE-04
\subsection*{Challenge VI.38.C4 solution}
\begin{solution}
Let
\[
f:X\to\mathbb P^n_Y
\]
be a \(Y\)-morphism. Its graph is
\[
\Gamma_f:X\longrightarrow X\times_Y\mathbb P^n_Y,
\qquad
x\longmapsto(x,f(x)).
\]
Let
\[
\Delta:\mathbb P^n_Y\longrightarrow
\mathbb P^n_Y\times_Y\mathbb P^n_Y
\]
be the diagonal. Relative projective space is separated over \(Y\), so \(\Delta\)
is a closed immersion. The graph square
\[
\begin{array}{ccc}
X & \xrightarrow{\ \Gamma_f\ } & X\times_Y\mathbb P^n_Y\\
\downarrow f && \downarrow {f\times\mathrm{id}}\\
\mathbb P^n_Y & \xrightarrow{\ \Delta\ } &
\mathbb P^n_Y\times_Y\mathbb P^n_Y
\end{array}
\]
is Cartesian. Therefore \(\Gamma_f\) is the base change of a closed immersion and
is itself a closed immersion.

Now suppose in addition that \(X\to Y\) is projective, with a chosen closed
immersion
\[
i:X\hookrightarrow\mathbb P^m_Y.
\]
The product of \(i\) with the identity gives a closed immersion
\[
X\times_Y\mathbb P^n_Y
\hookrightarrow
\mathbb P^m_Y\times_Y\mathbb P^n_Y.
\]
Composing with \(\Gamma_f\) yields
\[
(i,f):X\hookrightarrow
\mathbb P^m_Y\times_Y\mathbb P^n_Y.
\]
Finally the relative Segre embedding gives
\[
\mathbb P^m_Y\times_Y\mathbb P^n_Y
\hookrightarrow
\mathbb P^{(m+1)(n+1)-1}_Y.
\]
Thus
\[
\boxed{(i,f):X\hookrightarrow
\mathbb P^{(m+1)(n+1)-1}_Y}
\]
after Segre.  Graph embeddings therefore allow an existing projective embedding of
\(X\to Y\) to be enlarged so that a specified morphism to projective space is also
recorded by the projective coordinates.
\end{solution}

% VI38-DETAILED-CHALLENGE-05
\subsection*{Challenge VI.38.C5 solution}
\begin{solution}
Let
\[
S=k[x_0,\ldots,x_n],
\qquad
X\hookrightarrow\mathbb P^n_k
\]
be a closed subscheme, and let \(I\subseteq S\) be a homogeneous ideal defining
\(X\).  On a standard chart \(D_+(x_i)\), the ideal sheaf of the closed immersion is
represented by
\[
I_{(x_i)}=(IS_{x_i})_0\subseteq S_{(x_i)}.
\]
If
\[
I^{\mathrm{sat}}=I:S_+^\infty,
\]
then localization at \(x_i\) kills precisely the \(S_+\)-torsion relevant to
saturation, and
\[
IS_{x_i}=I^{\mathrm{sat}}S_{x_i}.
\]
Taking degree zero gives
\[
I_{(x_i)}=I^{\mathrm{sat}}_{(x_i)}.
\]
Thus the actual projective ideal sheaves agree:
\[
\boxed{\widetilde I=\widetilde{I^{\mathrm{sat}}}=\mathcal I_X}.
\]
This statement is already available from the standard-chart theory: projective
space sees the saturated homogeneous ideal, not irrelevant-ideal torsion.

There is also a global-sections formulation. For a homogeneous element \(F\in S_d\),
the condition that \(F\) define a section of \(\mathcal I_X(d)\) on every standard
chart is exactly that its local fractions lie in all the ideals
\(I_{(x_i)}\). Clearing powers of the \(x_i\) shows that this is the saturation
condition. Accordingly the saturated homogeneous ideal is recovered degree by
degree as
\[
\boxed{
I_X^{\mathrm{sat}}
=
\bigoplus_{d\ge0}
H^0\!\left(\mathbb P^n,\mathcal I_X(d)\right)
}
\]
under the standard identification
\(H^0(\mathbb P^n,\OO(d))\cong S_d\) for \(d\ge0\).

What belongs to later theory is not the chartwise saturation phenomenon itself, but
its systematic sheaf-theoretic framework and stronger consequences.  In later
chapters one proves, for coherent sheaves and finitely generated graded modules,
the precise correspondence modulo irrelevant torsion, finite generation of section
modules, Serre vanishing, eventual global generation, and cohomological criteria
for ampleness.  Those results allow one to pass back and forth between arbitrary
coherent sheaves and graded modules uniformly.  VI/38 needs only the special ideal
sheaf calculation above and should not import those later cohomological theorems.
\end{solution}
